


Written and committed BEFORE any run of Blocks P, Q or A. Branch \texttt{release}, started from tag \texttt{journal-v1} (\texttt{journal-v1.1} does not exist). Software only; \texttt{data/hardware} is not read. Result keys start with \texttt{r.} (repository), \texttt{ph.} (physics generator), \texttt{q.} (reruns) and \texttt{pa.} (protocol algorithms); no tagged result and no existing key is changed. There is no V6 preregistration; the number follows the release prompt.

\subsection*{0. Rules for every block}

\begin{itemize}
\item \textbf{Labels.} Every number is MEASURED\_SOFTWARE, SIMULATED, ANALYTICAL or NOT\_RUN. A negative result is an acceptable outcome of every task. Nothing is invented or adjusted.
\item \textbf{Synthetic data} is used for TRAINING and for clearly labelled stress tests only. Validation, calibration, thresholds and every real-data test metric use REAL data only.
\item \textbf{Seeds.} 42 for the headline run; seeds 0 to 4 for repeated model training; 30 fresh simulator seeds per cell in Blocks Q3 and A (design grid seeds 16000 + h, held-out grid seeds 18000 + h, Q3 grid seeds 19000 + h, h = 0 to 29). None of these seeds was used in an earlier block (earlier blocks used seeds below 10000).
\item \textbf{Splits.} NHTSA side: the sealed split of \texttt{docs/05\_data\_code/SEALED\_SPLIT.md} (development, validation, sealed fold) and the grouped, mode-stratified five-fold nested cross-validation of block F (the main estimate). VZCrash side: the supplied split with the de-duplicated test set primary and the strict set as sensitivity. Split by vehicle group; never by random window.
\item \textbf{Holm correction} over the family of the 20 hypotheses below. A hypothesis that is NOT\_RUN enters with p = 1 and is never supported. The p-value of each hypothesis follows the rule of \texttt{docs/02\_preregistration/PREREGISTRATION\_V5.md} section 0 (one-sided grouped bootstrap, or the cell-quantile rule for simulator grids). Anything defined after the results are seen is EXPLORATORY and is labelled so.
\item \textbf{Criteria are reused, not relaxed.} Each Q hypothesis keeps the numeric criterion of the earlier hypothesis it reruns (definitions in \texttt{PREREGISTRATION\_V3.md} for the X studies, \texttt{PREREGISTRATION\_V4.md} for the P studies, \texttt{PREREGISTRATION\_V5.md} for F and G). Where a Q item says "original", the original text governs.
\item \textbf{Wording.} The banned words of \texttt{CLAUDE.md} rule 7 are not used. Built-but-unmeasured items are called "built, not yet measured".
\item \textbf{Deviations.} Any change after a run has started is a dated, labelled deviation in \texttt{docs/07\_log/DECISIONS.md}. Items that cannot be completed in the session become NOT\_RUN with the reason in \texttt{results.json} and in \texttt{docs/07\_log/NOT\_RUN.md}.
\item \textbf{Gates before commit.} \texttt{check\_results} 0 errors, tests, ruff, link, wording and spelling checks, and the hash comparison of \texttt{results.json} and the tagged result files (Block R already passed it).
\end{itemize}

\subsection*{1. Block P: physics-based synthetic pulses}

\subsubsection*{1.1 Purpose and limit}
A generator whose pulses come from simple physical models with stated parameter ranges, used for training and for stress tests of sensor range. No claim about real crashes. Sentences for the paper (P5): abstract and methods: "Synthetic pulses were generated from stated physical assumptions and validated against real data on the checks in Section X." Every caption: "synthetic, generated from stated physical assumptions (docs/05\_data\_code/SYNTHETIC\_ASSUMPTIONS.md); validated against real data on the checks in the integrity report". The word "unbiased" is never used.

\subsubsection*{1.2 Models (primary generator PH-A)}
All equations, parameter ranges and sources are numbered in \texttt{docs/05\_data\_code/SYNTHETIC\_ASSUMPTIONS.md} (section "physics generator", items P-1 onward) with value and source or "assumed".
\begin{itemize}
\item \textbf{Frontal barrier:} body mass m, front structure as a Kelvin element (spring plus damper) with a nonlinear stiffness: linear rise to a plateau force, plateau, then bottoming out beyond a crush limit; impact speed v0; unloading with restitution e. The passenger-compartment deceleration is the output.
\item \textbf{Offset and other frontal:} the same model with a reduced engaged stiffness (overlap factor) and therefore longer duration.
\item \textbf{Side and pole:} narrow-contact model with a short door or pillar stage of low stiffness followed by a stiff stage; shorter duration and higher peak.
\item \textbf{Superimposed vibration and sensor mount:} damped sinusoids (frequencies and amplitudes assumed, stated) and a second-order low-pass mount response (cut-off assumed, stated).
\item \textbf{Parameter priors:} (a) quantities measured on the DEVELOPMENT fold of the NHTSA pool only (impact speed, reference delta-V, peak, duration, asymmetry) by moment matching of log-descriptors per mode; (b) published ranges only if the source was opened (none is relied on unless listed in the assumptions file with its URL and date opened). Vehicle mass is not in the records: it is an assumed prior, stated. The Kmec et al. (2021) numbers are not used.
\end{itemize}

\subsubsection*{1.3 Physical consistency checks (P2), per pulse}
Tolerances are fixed here: (i) integral of acceleration equals delta-V within 5\%; (ii) delta-V equals v0 (1 + e) within 5\%; (iii) absorbed energy equals half m v0 squared times (1 minus e squared) within 5\%; (iv) crush does not exceed the stated crush limit; (v) no negative closing velocity before separation; (vi) peak, duration and delta-V lie inside the real joint range of the mode (Mahalanobis distance of the log-descriptors from the development-fold mean below the 99.5\% chi-square limit for three degrees of freedom). The share failing each check is reported; failing pulses are dropped and the share dropped is stated.

\subsubsection*{1.4 Integrity gates (P3) that decide whether the generator may be used}
\begin{itemize}
\item \textbf{Provenance.} \texttt{SYNTHETIC\_PROVENANCE.md} (folder \texttt{docs/05\_data\_code}, written in Block P) lists the SHA-256 of the sorted record identifiers of the NHTSA development fold and of the VZCrash training split that each generator was fitted to. A generator may use only these. The earlier block-C/G generator was fitted to all 913 pooled records (\texttt{results/synthetic/generator\_params\_all913.json}), which includes the validation and sealed folds: it is REFITTED on the development fold only and the original and refit results are reported side by side (H-C3 stress test and the block G class pairs).
\item \textbf{Realism.} A real-versus-synthetic discriminator per mode: gradient-boosted trees (scikit-learn \texttt{HistGradientBoostingClassifier}, default settings, seed 42), trained on real development pulses against an equal number of synthetic pulses and scored on the REAL VALIDATION fold of the mode against freshly generated synthetic pulses (five times the real count). Two feature sets: D1, the pulse descriptors at the source rate (log peak, log duration, log delta-V, asymmetry, shape ratio, relative noise); D2, the 100 Hz, +-64 g estimator feature vector of \texttt{pp\_feat} plus a 24-point amplitude-normalised envelope. AUC with a 95\% bootstrap interval (2,000 resamples of the held-out sets). A mode PASSES when the point AUC of both D1 and D2 is at most 0.60; a mode whose upper interval end is above 0.60 is also named "uncertain" and reported separately. Failed modes are excluded and named. Also reported (no gate): energy-distance permutation test of the standardised log (peak, duration, delta-V) and the Frobenius norm of the difference of the correlation matrices, per mode. Pole and side failed earlier with the family generator; they are tested with the physics models and the outcome is reported either way.
\item \textbf{Independence from the algorithms.} The generator is frozen by hash BEFORE any algorithm is run on it: \texttt{results/synthetic/ph\_freeze.json} holds the SHA-256 of the generator source files and of the fitted-parameter file; it is committed before the earliest Q or A run. Every later change is listed in the change table of \texttt{SYNTHETIC\_PROVENANCE.md}; a change made after an algorithm result was seen is flagged ALGORITHM-INFORMED, and every result that used it is relabelled EXPLORATORY. Two structurally different alternative generators are fixed in advance: \textbf{FAM}, the earlier family generator (\texttt{fin\_synth}, Gaussian copula over half-sine, triangle and multi-peak shapes) refit on the development fold only; \textbf{PH-B}, a differently parameterised physics model (a series Maxwell-Kelvin element with an exponentially saturating force law, with priors placed on peak force and crush distance instead of stiffness and speed, and a different vibration set). Every conclusion of Blocks Q is checked on each alternative and reported as SAME, FLIPPED or GENERATOR-DEPENDENT (a verdict that holds on the primary generator and on only one alternative is GENERATOR-DEPENDENT). Each alternative passes the same gates 1 and 2 for its own modes.
\item \textbf{Settings that could favour a method} come from the real-data protocol and are listed with their source in the integrity report: sensor ranges 16, 32 and 64 g and the 100 Hz rate (\texttt{pp\_feat}, the sensor-selection study), the clip level (the range), the window (\texttt{config.py}: 1.0 s before and 1.8 s after the trigger for VZCrash; the NHTSA cells of \texttt{pp\_feat}), the 2 g trigger (\texttt{config.TRIGGER\_G}), the class boundaries 24, 40 and 56 km/h (\texttt{config.SEV\_DV\_KMH}).
\item \textbf{Negative controls.} Train on synthetic and test on real (TSTR); train on real and test on synthetic (diagnostic only, never a result); shuffled-label control (labels permuted within the training set; the metric must fall to the chance level). All three are reported per generator.
\item \textbf{Coverage statement.} Which real modes the generator represents (modes that passed) and the share of the 913 real records in those modes.
\end{itemize}

\subsubsection*{1.5 VZCrash-style surrogate (P4)}
The summary-statistics-only surrogate of the finishing run (\texttt{fin\_surr}) is refit on the VZCrash TRAINING split only and gated: (a) discriminator on VZCrash validation windows against surrogate windows (gradient-boosted trees on the window summary features and a 32-bin spectral envelope; point AUC at most 0.60); (b) nearest-neighbour no-copy check: the smallest distance from each surrogate window to the real training windows (standardised summary features) must not fall below the 1st percentile of real-to-real nearest-neighbour distances in more than 1\% of surrogate windows; (c) marginals: Kolmogorov-Smirnov distance at most 0.20 for each validated marginal; (d) spectra: median band-energy shares within 20\% relative of the real ones. Only fitted parameters are released (CC BY-NC 4.0, attribution to the VZCrash authors). It matches only the validated marginals; it is not called unbiased.

\subsection*{2. Block Q: reruns of every earlier algorithm}

\subsubsection*{2.1 Common protocol}
Training data may include the physics synthetic pulses (NHTSA side) and the surrogate windows (VZCrash side) when the generator passed its gates. Mixing ratio r (synthetic per real training row) is chosen on the REAL validation fold from \{0, 0.5, 1, 2, 5\}; r = 0 is the control with identical code. Evaluation is on REAL data only. The sealed NHTSA fold and the supplied VZCrash test set were already scored once in earlier blocks; each revised method is scored on them once more, labelled "revised after the negative", and the count of sealed-fold touches is reported. Grouped nested cross-validation is the main estimate on the NHTSA side. The primary generator is PH-A; FAM and PH-B are the alternatives of 1.4 (3). Original thresholds and criteria are reused.

\subsubsection*{2.2 Hypotheses (original criterion governs; what counts against is the original text)}
{\scriptsize\raggedright
\begin{longtable}{p{0.317\linewidth}p{0.317\linewidth}p{0.317\linewidth}}
\toprule
Id & Reruns & Planned evidence \\
\midrule
\endhead
H-Q1a & H-P1 direct classifier and augmented 1D CNN & nested CV and sealed fold; three generators \\
H-Q1b & H-P2 anchored interval (q calibrated on real validation only) & as above \\
H-Q1c & H-P3a abstention & as above \\
H-Q1d & H-X5 rank preservation & on the 176 original tests with synthetic rows inside the folds, if the split machinery of \texttt{x\_common} can take them; otherwise NOT\_RUN \\
H-Q1e & H-X6 anchored estimator, all four splits & as above \\
H-Q1f & H-X7 two-part fusion, original preregistered pairs & as above \\
H-Q2a & H-P4a compact forest (the int8 CNN P4d is exploratory) & supplied split, de-duplicated test primary, strict set as sensitivity \\
H-Q2b & H-P4b augmentation at the 1\% false-alarm target & as above \\
H-Q2c & H-P4c calibration variants & as above \\
H-Q2d & H-X1 16-bit forest & as above \\
H-Q2e & H-X2 trigger cascade & NOT\_RUN in advance: the cascade rule uses whole 16 s records and surrogate windows are trigger-aligned windows (the same reason as H-F2c) \\
H-Q2f & H-X3 group-aware training & as above \\
H-Q2g & H-X8 sequential rule & as above, if the six length-specific forests can be trained in the session \\
H-Q2h & the early-exit rule (N19) & decision times with surrogate-trained forests \\
H-Q3a & N28 equal-escalation comparison (P3 against the lowered-threshold table) & simulator, held-out grid of block B2 (12 cells), 30 fresh seeds (19000 + h), improved severity estimates from Q1 \\
H-Q3b & H-B2 P1+ against P0 & same grid and seeds \\
H-Q3c & P3b abstain routing & same grid and seeds \\
\bottomrule
\end{longtable}}
Items marked "if" are completed only if they run inside the session; otherwise they are NOT\_RUN with the reason and a dated deviation. No criterion is changed to make an item runnable.

\subsubsection*{2.3 Q4 summary table}
One row per hypothesis: original verdict, real-only rerun (r = 0), best-ratio rerun with interval, the verdict on each alternative generator, and whether the verdict changed (SAME, FLIPPED, GENERATOR-DEPENDENT). If nothing changes, the table says so.

\subsection*{3. Block A: protocol-layer algorithms (SIMULATED and ANALYTICAL)}

\subsubsection*{3.1 Numbered assumptions (A-1 to A-9; the full text, value and source mark are in \texttt{SYNTHETIC\_ASSUMPTIONS.md}, section "protocol")}
A-1 vehicle-powered nodes keep the receiver on when not transmitting (assumed). A-2 raw LoRa to a gateway, not LoRaWAN-compliant; nodes decode neighbours' alert frames and the gateway ACK within range at the same spreading factor and channel (assumed; range from the block G path-loss model). A-3 ACKs are authenticated (truncated HMAC, 8 bytes) and carry an event identifier and a coarse location cell (assumed). A-4 a GPS time pulse gives each node a clock error below 5 ms (assumed; the GPS module is in the design; the module's own accuracy is better, the bound includes interrupt latency). A-5 capture threshold 6 dB and log-normal shadowing as in the layers study (block G). A-6 multi-vehicle crashes produce near-simultaneous critical alerts from several nodes within a 10 s window (as the burst of the earlier grids). A-7 the India 433 MHz duty limit is a parameter (1\% and 0.1\%), UNVERIFIED. A-8 downlink and uplink loss as in the grids. A-9 the 97-byte alert and 26-byte ACK of \texttt{PACKET\_FORMAT.md}. The report states what happens if each assumption is removed (assumption sensitivity, 3.5).

\subsubsection*{3.2 Candidates}
\begin{itemize}
\item \textbf{OS-DR} (overhear-and-suppress, deadline-aware retries), built on the fast path of P1+ (flash, hold, severity table). A node that has sent an alert listens. It stops retrying only on an authenticated gateway ACK for its event cell and event identifier; never on an unauthenticated frame and never on a frame sourced by a neighbour. Retry times come from an analytical schedule: with attempt interval T (alert airtime plus listen window), deadline D = 10 s and per-attempt success probability s, the node draws its next attempt after a geometric number of attempt intervals with per-interval probability q, and q is the value that maximises the probability of an acknowledged alert before D under the duty budget in a renewal model with N contenders; the closed-form first-order solution is q = min(1, 1/N) with q = 1 inside the last two attempt intervals before D. N is a design parameter chosen on the design grid from \{1, 2, 4, 8, 16, 32\}. The schedule is a re-derivation in the style of the Markov models of confirmed-uplink delay (Khan et al., DCOSS 2019; known to this session only from the release prompt, to be opened in Block L, and no number from it is used) and is labelled a re-implementation of the style, not of their model. It is checked against simulation (3.5).
\item \textbf{OS-DR-N}: as OS-DR, and the retry probability is divided by (1 + 0.5 k), where k is the number of distinct neighbour alert frames for the same cell overheard in the last 10 s (a redundancy reduction only). It never delays or suppresses a node's initial transmission.
\item \textbf{GS-CA} (GPS-slotted, cluster-aware critical alerts): a critical alert is sent in slot h(vehicle id, event cell, epoch) mod K of the next GPS epoch, with a fresh epoch and hash for every retry; no beacon airtime is spent; slot length is the alert airtime plus a 20 ms guard (four times the clock bound); K is chosen on the design grid from \{1, 2, 4, 8, 16\} subject to K x slot length at most 8 s. Statistically this is slotted ALOHA with a public hash; the report says so if the results are the same.
\item \textbf{Baselines:} the best earlier policy (chosen on the DESIGN grid, see 3.3); slotted ALOHA with uniform slot choice on the same slot grid; an optimised slot-probability baseline re-implemented from the description of Tamang et al. (MSWiM 2023) given in the release prompt, which this session has not yet opened (one packet per node per slot frame, slot probabilities chosen by numeric optimisation to maximise the probability that at least one packet arrives for the design-grid cluster size); labelled a re-implementation, not the authors' code.
\end{itemize}

\subsubsection*{3.3 Grids}
Policies, simulator and metrics are those of block G (\texttt{fin\_g}, \texttt{world\_sim}), with the new features switched on through the world dictionary (all off by default; a test shows earlier results are reproduced when off). DESIGN grid: n in (75, 150), p in (0.1, 0.3), SF in (7, 9), duty in (1\%, 0.1\%), burst size b in (1, 20), seeds 16000 + h. HELD-OUT grid: n in (100, 200), p in (0.2, 0.4), SF in (8, 10), duty in (1\%, 0.1\%), b in (1, 30), seeds 18000 + h; 32 cells; 30 seeds per cell. "Burst cells" are the 16 held-out cells with b = 30. The best earlier policy is the earlier policy with the highest mean critical reach within 10 s on the DESIGN grid with no duty violation (ties: the lower node airtime share); it is fixed before the held-out run. All design choices (N, K, the OS-DR-N factor is fixed) use the DESIGN grid only. Nothing is tuned on the held-out grid and the simulator is not changed to produce a win.

\subsubsection*{3.4 Hypotheses}
\begin{itemize}
\item \textbf{H-A1.} OS-DR raises critical reach within 10 s (\texttt{reach10\_critical}, own-ACK reach, as in every earlier block) by at least 5 points over the best earlier policy in at least 60\% of the held-out cells (at least 20 of 32; interval of the paired seed difference above zero in each counted cell), loses by more than 2 points in no cell, and does not raise duty-cycle use (\texttt{node\_air\_share}) or non-severe escalation (\texttt{escalation}) above that policy's (paired interval upper end at most 0 in the cell mean). \textbf{Counts against:} fewer than 60\% of cells; any cell lost by more than 2 points; duty or escalation higher. Secondary, reported but not used for the verdict: gateway-reception reach (\texttt{reach10\_rx\_critical}) and an event-level reach (at least one alert of the event acknowledged within 10 s).
\item \textbf{H-A2.} GS-CA raises critical reach within 10 s by at least 5 points over both OS-DR and the Tamang-style baseline in at least 60\% of the burst cells (at least 10 of 16), with the same loss criterion. \textbf{Counts against:} fewer than 60\% of burst cells; any cell lost by more than 2 points.
\item \textbf{H-A3 (security regression).} The 48 attack cases and the bounded model check, extended with an attacker that forges ACKs or neighbour frames to suppress alerts, show no new counterexample for OS-DR. OS-DR-N is dropped if a spoofed neighbour frame can delay an initial transmission. \textbf{Counts against:} any new counterexample.
\end{itemize}

\subsubsection*{3.5 Reporting (A4)}
Regime map update (when each policy wins); assumption sensitivity (each of A-1 to A-9 varied or removed; which conclusions flip); analytical-versus-simulated agreement of the retry schedule in a controlled scenario (N contenders, fixed success probability, deadline D; reported as the difference between formula and simulation with an interval); duty violations, which must be 0. If neither candidate passes on the held-out grid, both are added to the negatives register.

\subsection*{4. Blocks L and E (process, no hypotheses)}
L: update the reference lists from sources that were opened; \texttt{LITERATURE\_STATUS.md} was named as supplied but is not in the repository, so the reference work starts from what the session can open itself (decision logged 2026-11-13); nothing is cited that was not opened; Scopus is not run by Claude. E: an independent claim audit of README, project status, claim lists and the paper against \texttt{results.json}; regenerate \texttt{paper/generated}; tag \texttt{release-v1}; build the release archive without raw third-party data; list the push and Zenodo commands for the author.

